\paragraph{The endpoint-invariant repair is masked by a strictly later LUT input.}
The previously unresolved SPMV masking mechanism is now identified from a separate, disabled-by-default native timing observer. It exports the existing Tatum graph, constraints, edge delays and arrival/required tags at 17 digits. The frozen repair implementation and its 45 source bindings remain unchanged. Three successful analysis-only observations preserve their original route and complete endpoint vectors exactly: observer off and on for the accepted plateau, and observer on for the rejected successor. No routing is repeated. Each exported graph contains 66118 nodes, 176948 edges and 382088 tags.

The successor changes only one timing-edge delay, corresponding to 6997/pin~1, and only that connection's receiving node changes setup arrival. Its next LUT output is still determined by 6927/pin~1. The repaired input was already 129.686484 ps earlier than that unique winning input; after repair it is 315.512061 ps earlier. All incoming arrival-plus-edge maxima reproduce native float32 arithmetic in both states. Thus this is strictly dominant-input masking, not an equal-delay tie or a failure to adopt the physical shortcut. All 129 reachable downstream endpoints retain their arrival and required tags. The candidate remains legal but both acceptance policies reject its unchanged CPD and endpoint vector. The accepted plateau remains at CPD 7.034697713237392 ns, one SB pair, 967 CB sources and wirelength 148445. The earlier two-component setup gain does not unlock subsequent CPD improvement in this sequence.

\paragraph{Current endpoint sensitivity changes where the next search is spent.}
A new read-only selector starts from the actual minimum-slack endpoint and recursively retains all exact maximum-arrival predecessors for its launch domain, recording unresolved recurrences as unknown. This is distinct from downstream reachability and from selecting every connection with a near-critical pin slack. The SPMV current worst endpoint is \texttt{fetcher.fifo\_col\_id\_30.gb.D[0]}. Its complete traced cone has 14 nodes and 13 edges, with six external routed connections: 6157/pin~12, 6373/pin~1, 6923/pin~1, 6927/pin~1, 6998/pin~1 and 7773/pin~1. There are no unresolved non-source nodes.

All six current physical connection delays equal their exact full-device SOURCE-only shortest distances in the actual zero-RC graph, even with background occupancy and retained branches relaxed. Returned path edges are checked against the XML-bound graph and native-rounded switch delays. The first three queries cover the dominant masking input's upstream cone; the full-endpoint pass reuses those exact cases and queries only the other three. This establishes no remaining physical-distance gain for those six fixed-terminal connections. It is a conditional limit of this route-only intervention, not a global placement, pin-assignment, logic or FPGA-routing optimum. Endpoint-Pareto work outside the worst cone is not excluded. The selector currently supplies scheduling evidence; it has not yet been promoted as an integrated repair-controller optimization or claimed to improve quality.

\paragraph{Diagnostic cost and preserved engineering failures.}
The three successful observation child processes take 47.320 s together, including 25.559 s of native execution. One earlier observation takes 11.985 s before its automatically selected TBB configuration hits the thread guard; the successful observer explicitly matches the baseline's serial engines. A parser error after two successful observations is fixed without repeating those observations. Initial missing-server dependency, recovered executable packaging and parser failures remain in the evidence. The three build attempts take 13.038, 194.893 and 162.088 s, reported as engineering cost rather than route cost. Causal graph analysis takes 8.179 s, automatic endpoint-cone selection 3.881 s, and the two source-query batches 4.460 and 4.695 s, with three query cases reused in the latter. These are development measurements with nested times kept separate; no U-version historical timing is substituted.

The previously prepared complete physical-audit cache is now exercised on fresh native observations: independent audit components take 4.343 and 4.394 s, inside complete evaluation times of 17.619 and 16.456 s. This establishes successful use, not a matched whole-process speedup. Further research moves to untouched BNN Tf, strong native A0 and old R1 routes with fresh endpoint selection and independent strict/endpoint branches. BNN has a different physical graph and must receive its own graph binding before capacity queries. These registered observations are not yet presented as repair results or a terminating experimental matrix.
